\subsection{Problem 14. Volumes by disks and shells}
A solid is obtained by rotating the region bounded by $y=x^{2}$, $y=0$, and $x=1$ about:
\begin{enumerate}[label=(\alph*)]
    \item The $x$-axis (use the disk/washer method).
    \item The line $x=2$ (use the shell method).
\end{enumerate}

\begin{enumerate}[label=(\alph*), start=3]
    \item Set up (but do not yet evaluate) the integral for the volume in case (a), then compute it.
    \item Set up and evaluate the integral for the volume in case (b). Simplify as much as possible.
    \item Compare the two methods (disks vs shells) in each case and explain which is more convenient and why. Would it still be the case if the region were more complicated?
\end{enumerate}

\subsection*{Strategy}
\subsection*{Solution}
\begin{center}
\begin{tikzpicture}
    \begin{axis}[
        axis lines=middle,
        xlabel=$x$, ylabel=$y$,
        xmin=-0.5, xmax=2.5,
        ymin=-0.5, ymax=1.5,
        width=0.6\textwidth,
        height=0.6\textwidth
    ]
        % Plot y = x^2
        \addplot[name path=A, blue, thick, domain=0:1] {x^2};
        
        % Plot y = 0 (x-axis)
        \addplot[name path=B, black, thick, domain=0:1] {0};
        
        % Plot x = 1
        \addplot[black, thick] coordinates {(1,0) (1,1)};

        % Shade the region
        \addplot[gray!30] fill between[of=A and B];
        \node at (axis cs:0.7, 0.2) {Region};

        % Draw Axis of Rotation for (b)
        \addplot[red, dashed, thick] coordinates {(2,-0.5) (2,1.5)};
        \node[red, right] at (axis cs:2, 0.8) {Axis $x=2$};
        
    \end{axis}
\end{tikzpicture}
\end{center}
\textbf{(c) Volume about the x-axis (Disk Method)}
Since we are rotating about the x-axis, the cross-section at $x$ is a disk with radius $R(x) = x^2$.
\[ V_a = \int_{0}^{1} \pi [R(x)]^2 dx = \pi \int_{0}^{1} (x^2)^2 dx \]
Computing the integral:
\[ V_a = \pi \int_{0}^{1} x^4 dx = \pi \left[ \frac{x^5}{5} \right]_{0}^{1} = \pi \left( \frac{1}{5} - 0 \right) = \boxed{\frac{\pi}{5}} \]

\textbf{(d) Volume about the line x=2 (Shell Method)}
We use cylindrical shells with thickness $dx$.
\begin{itemize}
    \item \textbf{Radius ($r$):} The distance from the slice at $x$ to the axis $x=2$ is $r = 2 - x$.
    \item \textbf{Height ($h$):} The height of the function at $x$ is $h = x^2$.
\end{itemize}
The volume formula is $V = \int_{a}^{b} 2\pi r h \, dx$.
\[ V_b = \int_{0}^{1} 2\pi (2-x)(x^2) dx \]
Evaluating the integral:
\[ V_b = 2\pi \int_{0}^{1} (2x^2 - x^3) dx \]
\[ V_b = 2\pi \left[ \frac{2x^3}{3} - \frac{x^4}{4} \right]_{0}^{1} \]
\[ V_b = 2\pi \left( \left(\frac{2}{3} - \frac{1}{4}\right) - 0 \right) = 2\pi \left( \frac{8-3}{12} \right) = 2\pi \left(\frac{5}{12}\right) = \boxed{\frac{5\pi}{6}} \]

\textbf{(e) Comparison of Methods}
\begin{itemize}
    \item \textbf{Case (a):} The Disk method was very convenient because the rotation axis ($y=0$) was a boundary of the region, and the function $y=x^2$ was explicitly given. Using shells here would require integrating with respect to $y$ (horizontal shells), involving $x=\sqrt{y}$, which is slightly more complex.
    \item \textbf{Case (b):} The Shell method was convenient because we could continue using the variable $x$. If we used the Washer method (slices perpendicular to $x=2$, i.e., horizontal $dy$), we would need:
    \begin{itemize}
        \item Outer radius $R_{out} = 2 - 0 = 2$ (distance from axis to y-axis).
        \item Inner radius $R_{in} = 2 - \sqrt{y}$ (distance from axis to curve).
    \end{itemize}
    The integral would be $\pi \int_0^1 (2^2 - (2-\sqrt{y})^2) dy$, which involves more algebra to expand.
    \item \textbf{Conclusion:} The shell method is generally better when rotating about vertical lines for functions given as $y=f(x)$, as it avoids solving for $x$ and dealing with potentially complicated inverse functions.
\end{itemize}
\pagebreak
